In rescue-service incident response, command belongs to a person: the incident commander, appointed by the organization. When that commander is cut off from the incident information system, the organization may hand command to a deputy or a superior, and the system must then decide which recorded decisions are authoritative. Clocks can order one person's decisions, but they cannot say whose authority was valid when two people each believed they were in command. We separate a person's right to command from a node's right to write. Authority comes from signed role grants and a succession rule that the organization fixes in advance; nodes exercise that authority but hold none of their own, and the network never decides who is in command. A decision is binding when it is admitted to the authoritative journal under a valid grant. Decisions made under a grant that had already been superseded are kept and labelled as residual, never merged or dropped. While the commander is cut off, a design can give at most two of three things: the commander keeps recording, command moves without the commander, and only one disconnected edge accepts decisions. We therefore treat succession as a choice of policy, compare six succession policies by what each one gives up, and place Authority-Aligned Sequencing (\systemname{}), our prototype, as one point in that space: immediate promotion through the regional core with a lease. We report what the prototype does and does not guarantee today, using a two-edge partition experiment, TLA+ model checking, and baselines against Raft and a last-writer-wins CRDT. \pending{A policy comparison across three partition cases is planned; no results yet.}
